\section{Sim2Real、Real2Sim 与合成数据}

本章先消化高频术语，因为后面的产业判断都建立在这些概念上。Sim2Real 是从仿真走向现实，Real2Sim 是把现实映射进仿真，Synthetic Data 是通过仿真、算法和程序化流程生成可训练数据。三者不是并列 buzzword，而是同一个工程闭环的不同位置：真实世界给出问题，仿真扩充问题，模型在合成数据上训练，部署结果再反向校准仿真。

谢晨在 Cruise 的经验提供了一个关键教训。早期团队用游戏工程师复刻旧金山，看起来像 Matrix，但感知部门无法直接使用这些数据，因为它们没有被机器学习目标定义。后来真正有效的做法，是先建立评价系统，判断合成数据是否能提升模型，再反过来改造资产、场景、传感器噪声和数据生成链路。

\begin{knowledgebox}{术语消化：仿真、合成数据和迁移}
\begin{tabularx}{\textwidth}{p{0.20\textwidth}p{0.35\textwidth}X}
\toprule
术语 & 第一层含义 & 在本期中的关键问题 \\
\midrule
Sim2Real & 在 simulation 中训练或生成数据，再迁移到真实世界 & 真实部署是否成功，决定仿真是否有用。 \\
Sim2Real gap & 仿真与真实之间的分布差异 & 不是追求零差距，而是判断差距小到什么量级时数据开始有用。 \\
Synthetic Data & 由仿真、算法、程序化流程或人机协作生成的数据 & 关键不是生成数量，而是信息密度、质量和模型效用。 \\
Real2Sim & 把真实场景、资产、物理参数和机器人本体映射进仿真 & 视觉重建相对容易，物理参数和交互力学才是难点。 \\
Evaluation Loop & 用模型训练和真实部署结果反推数据质量 & 没有评价系统，数据团队很难知道该改哪里。 \\
\bottomrule
\end{tabularx}
\end{knowledgebox}

\subsection{从漂亮玩具到机器学习数据}

本节回答一个具体问题：为什么“看起来真实”的仿真不等于“对模型有用”的仿真。Cruise 早期的仿真系统能复刻街区，但如果传感器噪声、光照分布、标注、场景概率和长尾事件不符合模型训练需要，感知系统就无法受益。谢晨的第一步不是加画质，而是建立评价系统。

评价系统有两类指标。第一类是绝对真实性：光线、色彩、传感器噪声、几何、标注、场景分布和真实数据是否对齐。第二类是效用性：把这批合成数据喂给模型之后，模型在目标任务上的表现是否提升。讲者用浓咖啡和兑水冰咖啡作比喻：后续数据越来越看重“每单位数据给模型带来多少信息增量”，而不只是体积。

\begin{importantbox}{评价系统优先}
合成数据团队最早要问的不是“怎么生成更多”，而是“什么数据能让模型变好”。没有评价系统，仿真团队只能优化视觉观感；有了评价系统，团队才能知道要改传感器噪声、场景分布、资产物理参数，还是数据质检链路。
\end{importantbox}

可以把数据效用写成一个简化目标：
\[
U(D_s)=\Delta M_{\mathrm{target}}-\lambda C(D_s)-\gamma G(D_s,D_r)
\]
其中，$D_s$ 表示合成数据，$D_r$ 表示真实数据，$\Delta M_{\mathrm{target}}$ 表示目标模型指标提升，$C(D_s)$ 表示生成与质检成本，$G(D_s,D_r)$ 表示合成和真实之间的关键 gap，$\lambda$ 和 $\gamma$ 是成本与 gap 的权重。这个公式不是访谈原文，而是对访谈方法论的教学化表达：有效合成数据必须同时考虑模型收益、生产成本和真实差距。

\subsection{合成数据流程与配比}

接下来本节把流程落到一个例子：路上行人突然窜出。这类 corner case 在真实数据中少，但对安全很关键。合成数据流程不是把一个视频复制一万份，而是围绕一个真实问题做系统性变体：不同年龄、身形、衣服、路线、车道、天气、光照、交通状态和传感器噪声。然后让算法在仿真里跑一万遍，发现感知、预测和规划是否还会出错。

配比没有固定答案。对自动驾驶整体盘子而言，真实车队每天产生大量真实数据，合成数据可能在整体占比里不是最高；但对稀缺 corner case，真实样本可能一年也不够，因此一个真实问题会配上成千上万个合成变体。节目中给出的经验是：整体上曾经大约 30\% 使用合成数据，而在长尾场景上可以超过 1:99 或 1:100 的真实/合成比例。

\begin{knowledgebox}{流程消化：从一个问题到一批训练数据}
\begin{tabularx}{\textwidth}{p{0.20\textwidth}p{0.35\textwidth}X}
\toprule
步骤 & 做什么 & 为什么重要 \\
\midrule
真实问题定位 & 找到一次危险或失败案例 & 数据生成必须从真实风险出发。 \\
场景参数化 & 改行人、车道、天气、光照、交通流 & 让模型看到同一问题的多种表现。 \\
传感器仿真 & 模拟相机、雷达、噪声和观测视角 & 模型最终看到的是传感器数据。 \\
自动标注与质检 & 生成 2D/3D/4D 标签并检查质量 & 标注错误会直接污染训练。 \\
模型回灌 & 用数据训练或测试感知/预测/规划 & 效用要由模型指标和部署结果证明。 \\
\bottomrule
\end{tabularx}
\end{knowledgebox}

\begin{warningbox}{配比不能脱离任务讨论}
说“真实数据 1\%，合成数据 99\%”容易误导。正确问法是：对哪个任务、哪个失败模式、哪个模型模块、哪个部署阶段，合成数据的边际效用更高？整体数据盘子、长尾场景和强化学习环境的配比可能完全不同。
\end{warningbox}

\subsection{本章小结}

Sim2Real 和 Real2Sim 构成数据闭环。合成数据的价值不由画面真实感决定，而由两类评价决定：它和真实世界关键变量是否对齐，以及它是否能让目标模型在真实任务上变好。

